\documentclass[aspectratio=169]{beamer}
\usepackage[utf8]{inputenc}
\usepackage{hyperref}
\usepackage{amsmath,amsfonts,amsthm,bm}
\usepackage{color}
\usepackage{minted}
\usepackage{graphicx} % Allows including images
\usepackage{subcaption}
\usepackage{booktabs} % Allows the use of \toprule, \midrule and \bottomrule in tables
\usepackage{tikz}
%\usepackage{pgfplots}
\usepackage{listings}
\usepackage{courier}
\usepackage[version=4]{mhchem}
\usepackage{array}
\setminted{fontsize=\scriptsize}
\lstset{ %
    basicstyle=\scriptsize\ttfamily, % fonts that are used for the code
    breakatwhitespace=false,         % sets if automatic breaks should only happen at whitespace
%breaklines=true,                 % sets automatic line breaking
%captionpos=b,                    % sets the caption-position to bottom
    commentstyle=\color{gray}\textit,    % comment style
    keepspaces=true,                 % keeps spaces in text, useful for keeping indentation of code (possibly needs columns=flexible)
    keywordstyle=\color{blue},       % keyword style
    language=Python,                 % the language of the code
%otherkeywords={*,...},          % if you want to add more keywords to the set
    rulecolor=\color{black},         % if not set, the frame-color may be changed on line-breaks within not-black text (e.g. comments (green here))
    showspaces=false,                % show spaces everywhere adding particular underscores; it overrides 'showstringspaces'
    showstringspaces=false,          % underline spaces within strings only
    showtabs=false,                  % show tabs within strings adding particular underscores
    stringstyle=\color{red}, % string literal style
    tabsize=4,                       % sets default tabsize to 2 spaces
    columns=fixed                    % Using fixed column width (for e.g. nice alignment)
}

\hypersetup{
    colorlinks=true,
    linkcolor=red,
    filecolor=magenta,
    urlcolor=red,
}

\DeclareMathOperator*{\argmax}{argmax}
\DeclareMathOperator*{\argmin}{argmin}
\let \vec \mathbf

\newcommand{\classname}{NANO266}
\newcommand{\classyear}{Fall 2024}
\mode<presentation> {
    \usetheme{CambridgeUS}
    \setbeamertemplate{footline}[text line]{%
        \parbox{\linewidth}{\vspace*{-8pt}\classname\hfill\classyear\hfill\insertpagenumber}}

    %\setbeamertemplate{footline}[page number]
    \setbeamertemplate{navigation symbols}{}
}


\title[\classname Lab 2 - A Study in Silicon]{\classname~- Quantum Mechanical Modeling of Materials and Nanostructures\\Lab 2 - A Study in Silicon}

\author{Shyue Ping Ong}
\institute[UCSD]{University of California, San Diego\\
\medskip
}
\date{\classyear} % Date, can be changed to a custom date

\begin{document}


    \begin{frame}
        \titlepage % Print the title page as the first slide
    \end{frame}


    \begin{frame}{The Material Powering Our Digital Age}
        \begin{figure}
            \centering
            \includegraphics[width=0.8\linewidth]{lectures/figures/Lab_2_Silicon_Valley.png}
        \end{figure}

    \end{frame}


    \begin{frame}{Aims of Lab 2}
        \Large{
            \begin{enumerate}
                \item Learn about convergence of QM calculations, specifically with respect to energy cutoffs and k-points.
                \item Efficient conduct of calculations using simple programming.
                \item Derivation of simple mechanical properties from DFT calculations
            \end{enumerate}
        }
    \end{frame}

    \begin{frame}{Quantum Espresso}
        \begin{figure}
            \centering
            \includegraphics[width=0.3\linewidth]{lectures/figures/Lab_2_Quantum_Espresso.png}
        \end{figure}
        Integrated suite of open-source computer codes for electronic-structure calculations and materials modeling at the nanoscale. It is based on density-functional theory, plane waves, and pseudopotentials. (http://www.quantum-espresso.org/)\newline
        \newline
        For this lab, we will specifically be using PWSCF. Documentation available \href{https://www.quantum-espresso.org/Doc/pw_user_guide/}{here}
        (you will definitely need this).
    \end{frame}

    \begin{frame}[fragile]{Templated PWSCF Input File}

        \begin{alertblock}{Parameters}
            \{\} denotes variables that will be modified by the programming script.
        \end{alertblock}

        \footnotesize
        \begin{verbatim}
&control                   # This is the control section
  calculation = 'scf' ,    # Specifies that we are doing a static SCF calculation.
  outdir = './tmp' ,
  pseudo_dir = './' ,      # Location of pseudopotential files.
  tprnfor = .True.,        # Specifies that you want forces to be printed. Used
  tstress = .True.,        # in subsequent questions.
/
&system                    # Specifies the structure
  ibrav = 2,               # For PWSCF, ibrav = 2 denotes an FCC cell.
  celldm(1) = {alat},      # This specifies the lattice parameter of the fcc cell.
  nat = 2,                 # We have two Si atoms per unit cell.
  ntyp = 1,                # There is only one type of atom (Si)
  ecutwfc = {ecut} ,       # This stipulates the energy cutoff.
        \end{verbatim}


    \end{frame}


    \begin{frame}[fragile]{Templated PWSCF Input File, contd.}
        \footnotesize
        \begin{verbatim}
&CELL
/
ATOMIC_SPECIES             # Specify the pseudopotential for each species.
  Si   28.055  {pseudopotential}
ATOMIC_POSITIONS crystal   # Specifies the atomic positions in frac. coords
  Si      0.00    0.00    0.00
  Si      0.25    0.25    0.25
K_POINTS automatic         # Specifies the k-point grid to be used
  {k} {k} {k}   0 0 0
        \end{verbatim}

    \end{frame}

    \begin{frame}[fragile]{Automation of Calculations with Python}

        \textbf{run\_pw.py} script has been provided.

        \inputminted{python}{code/Lab_2_code.py}

    \end{frame}

% \begin{frame}[allowframebreaks]{Bibliography}
%     \bibliographystyle{unsrt}
%     \bibliography{refs}
% \end{frame}

    \begin{frame}
        \Huge{\centerline{The End}}
    \end{frame}

\end{document}

